\chapter{Conclusion}
\label{ch:conclusion}

% NOTE: future work lives in Ch.~\ref{ch:limitations-future-work} (per the
% ADR-055 amendment), so this chapter stays short and does not repeat it.

This thesis began with a question forced by law rather than by technology: how reliably
can current open-weights document \acp{VLM}, deployable on hardware a German tax
firm already owns, convert business-to-business invoices into the legally grounded
\ac{EN}~16931 form --- and which stage of such a pipeline limits what is achievable? The
professional-secrecy regime of \S203~\ac{StGB} makes the local constraint an operational
premise, not a preference (\S\ref{sec:motivation}), and the thesis asked what is
honestly attainable within it.

\section*{The answer}
\addcontentsline{toc}{section}{The answer}

The answer has a shape, not just a number. On 39 real invoices the system never saw
during development, zero-shot header-field extraction reaches a mean per-invoice
\ac{F1} of 0.88 against an answer key adjudicated across three independent channels
(\S\ref{sec:results-heldout}). The error profile leans the right way for the domain:
precision 0.95 against recall 0.85, with four errors in five being omissions a reviewer
can see rather than fabrications a reviewer must catch --- the honest-null contract of
the design (\S\ref{sec:design-honest-null}) showing up in the measurements. The cost
structure is equally clear: language costs little, degraded capture costs more than
eleven points, and the synthetic-corpus figures are upper bounds for camera-captured
paper (\S\ref{sec:lim-corpus}).

The stage that limits the pipeline is reading. The structuring stage exceeds 0.97 on
perfect transcripts while the same stack, differing only in that the transcript
comes from a real reader, reaches 0.85 on the same corpus at matched precision
(\S\ref{sec:results-baseline}); and in the per-miss attribution --- taken with the
survey-era reader and labelled as the diagnostic it is --- roughly half of all
individual errors were values the transcript never contained
(\S\ref{sec:results-attribution}). Effort spent
below the binding constraint is wasted: the single cheapest improvement made in this
project was swapping the reading model --- one configuration line --- and the most
expensive intervention attempted, a pre-registered low-rank adaptation of the
structuring stage, regressed in all four cells of its design and specifically raised
the rate of values invented for absent fields (\S\ref{sec:results-adaptation}).
Both outcomes point the same way: at this scale, buy reading quality, do not tune
structuring.

\section*{The contribution to remember}
\addcontentsline{toc}{section}{The contribution to remember}

If one thing from this thesis outlives its model rankings --- and model rankings from
this period will not live long --- it should be the measurement discipline of
Chapter~\ref{ch:measurement-validity}. The scoring corrections documented there moved
reported figures by amounts comparable to the differences between models; one
un-repaired comparison rule inverted the ordering of two candidate readers; and one
precision-mismatched baseline, caught because a pre-registered rule required the
re-measurement, would have flipped the adaptation conclusion from harmful to harmless
(\S\ref{sec:validity-confound}). None of these defects announced themselves; every one
was found by treating the measuring apparatus as a system under test, with frozen
generations as the mechanism that makes an instrument repair measurable in isolation
from a model change. A reported score is a claim about a model \emph{and} about the
ruler it was measured with, and this thesis's experience suggests the ruler's share of
the uncertainty is routinely underestimated.

\section*{What this thesis does not claim}
\addcontentsline{toc}{section}{What this thesis does not claim}

The boundaries are stated as plainly as the results. The knowledge-graph and query
layers were designed and their hypotheses registered, but they were not built or
evaluated, and no finding here bears on them; the registered hypotheses are reported as
unevaluated, never quietly dropped (\S\ref{sec:unevaluated}). No cloud comparison
was run. The held-out corpus is small and single-annotator, so its figures establish a
level and a direction, not precise values; the structuring model was never
competitively selected, unlike the reader; and everything measured concerns German and
European \ac{B2B} invoices, nothing else. One boundary deserves naming here because it
concerns the deployment premise itself: the hardware claim is two-tier, and the tiers carry
different evidence. Every number reported above was produced on premises-class hardware ---
nothing needed more than a single 24\,\ac{GB} graphics processor, running open-weights models
under the author's control at every step --- while the laptop floor, whose memory ceiling caps
the reader at roughly half of native page resolution, has a measured structuring stage but an
unmeasured reading quality (\S\ref{sec:method-repro}). The delivered pipeline is open-weights,
makes no network call, and fits the floor --- but its reading quality at the floor was not
measured, and no figure here should be read as a floor figure
(\S\ref{sec:lim-hardware}). Within those boundaries, the claims are as
strong as the committed evidence behind them --- every number in this manuscript is
generated from the archived measurement artefacts, none transcribed by hand ---
and no stronger.

\section*{Outlook}
\addcontentsline{toc}{section}{Outlook}

The path forward is ordered by the evidence, not by ambition, and
Chapter~\ref{ch:limitations-future-work} states it in full: improve capture and
reading before touching structuring; treat the unbuilt layers as the next vertical,
carrying the same measurement discipline upward; and re-run the sealed harness as the
model generation turns over, which the frozen-artefact design makes cheap. The
practical conclusion for the profession that motivated this work is modest and
useful: a locally deployable pipeline that reads most of a real invoice correctly,
fails visibly rather than plausibly, and comes with an evaluation apparatus that can
say --- with evidence --- when it has become good enough to trust.
